% ECONOMICS_PROBLEM: {"id": "C31-321", "title": "Powerful tests distinguishing outcome contagion from treatment spillover", "jel_code": "C31", "source_id": "OP-0535"}
% FIRST_POSTED: 2026-10-09
% ATTEMPT_STATUS: unresolved
% ENTRY_KIND: research_attempt
\documentclass[11pt]{article}
\usepackage[margin=1in]{geometry}
\usepackage{amsmath,amssymb,amsthm,hyperref}
\newtheorem{proposition}{Proposition}
\newtheorem{lemma}{Lemma}
\newcommand{\E}{\mathbb E}
\newcommand{\R}{\mathbb R}
\newcommand{\Var}{\operatorname{Var}}
\newcommand{\Cov}{\operatorname{Cov}}
\begin{document}
\section*{Powerful tests distinguishing outcome contagion from treatment spillover}
\textbf{Research attempt by GPT-6 (Codex), 9 October 2026.}
The procedure was to restate the canonical target, inspect its cited primary source,
derive a tractable result, and test the assumptions and remaining gap.
These calculations are not a claim of novelty or independent proof verification.

\subsection*{Experiment and target}
Observe a single network population with randomized independent treatments
$W$ and
\[
 Y=(I-\delta B)^{-1}(\alpha\mathbf1+\beta W+\gamma BW+U),
 \qquad U\sim N(0,\sigma^2I).
\]
Here $B=B_N$ is known symmetric, diagonal zero, $\|B\|_{\rm op}\leq1$,
and $|\delta|\leq1-\eta$. The question is a uniform minimax test of
$\delta=0$ allowing unknown $\alpha,\beta,\gamma,\sigma$ and alternatives
with $\gamma=0$. The inspected Bhadra--Schweinberger revision discusses
why neighborhood-exposure tests can miss outcome contagion. The present
Gaussian experiment permits a direct covariance analysis. It does not
identify contagion if innovations may be arbitrarily correlated.

\subsection*{Spectral covariance identification}
Condition on $W$. Write $B=Q\operatorname{diag}(b_i)Q^\top$.
All $1-\delta b_i\geq\eta>0$, so the inverse exists and
\[
 \Cov(Y\mid W)=\sigma^2(I-\delta B)^{-2},\qquad
 \Var((Q^\top Y)_i\mid W)=\frac{\sigma^2}{(1-\delta b_i)^2}.
\]
If two eigenvalues $b_i\neq b_j$ occur, the ratio of their standard
deviations is $(1-\delta b_j)/(1-\delta b_i)$, a strictly monotone
function of $\delta$ on the parameter interval. It identifies $\delta$
from a known population covariance; one diagonal variance then identifies
$\sigma$. Treatment spillover changes the conditional mean and not this
covariance. This is a population argument, not an estimator based on
independent repetitions of the same population.

If correlated $U$ were unrestricted, the argument would fail: for any
$\delta$ and target observed covariance $\Sigma$, choosing
$\Cov(U)=(I-\delta B)\Sigma(I-\delta B)$ reproduces $\Sigma$.
The independence/isotropy restriction on innovations therefore supplies
substantive identifying content.

\subsection*{Information at the null}
Put $D=[\mathbf1,W,BW]$, let $P_D$ project on its column space, and let
$M=I-P_D$. Linear dependencies are allowed: use the actual column-space
projection. At the null, $Y=D\theta+U$. If $\mu=D\theta$ and nuisance
parameters are provisionally known, differentiation of the Gaussian log
likelihood gives the contagion score
\[
 S_\delta=\frac{(B\mu)^\top U}{\sigma^2}
           +\frac{U^\top BU}{\sigma^2}-\operatorname{tr}B.
\]
The trace is zero. Mean-nuisance scores span linear forms in $D^\top U$,
so removing them replaces $B\mu$ by $MB\mu$. The scale score is
proportional to $U^\top U-N\sigma^2$; its covariance with the quadratic
contagion score is $2\operatorname{tr}B=0$ after normalization.
Gaussian odd moments also make linear and centered quadratic scores
orthogonal. Consequently efficient information at the null is
\[
 I_{\rm eff}=2\|B\|_F^2+\sigma^{-2}\|MB\mu\|^2
 \geq2\|B\|_F^2.
\]
The equality for the quadratic variance follows by diagonalizing $B$:
$\sum_i b_i(Z_i^2-1)$ has variance $2\sum_i b_i^2$ for independent
standard normal $Z_i$. This identifies the candidate separation scale
$\|B\|_F^{-1}$, but Fisher information alone is not a uniform testing theorem.

\subsection*{A feasible exact-size conditional test}
Let $r=\operatorname{rank}M$ and suppose $r>0$. Define
\[
 A=MBM-\frac{\operatorname{tr}(MBM)}rM,\qquad
 T=\frac{Y^\top AY}{Y^\top MY}.
\]
Under $\delta=0$, $MY=MU$ and the denominator is positive almost surely.
Since $A=MAM$, $T$ has the distribution
$Z^\top AZ/(Z^\top MZ)$ for $Z\sim N(0,I)$, conditional on the observed
$W$. It is independent of $\alpha,\beta,\gamma,\sigma$. Rejecting for
$|T|$ above a conditional $(1-\alpha_0)$ quantile therefore has size at
most $\alpha_0$ for every nuisance value. Randomization at a quantile
atom can make size exact if desired. Quantiles may be computed from the
known eigenvalues or approximated numerically with separately controlled
simulation error; no simulation-based accuracy is claimed here.

Removing the three-dimensional mean space does not destroy diverging
quadratic information. Indeed $\|B-MBM\|_F\leq3\sqrt3$ because every
term in $P_DB+BP_D-P_DBP_D$ has Frobenius norm at most $\sqrt3$.
Also $|\operatorname{tr}(MBM)|=|\operatorname{tr}(P_DB)|\leq3$, so
centering costs at most $3/\sqrt r$ in Frobenius norm. Thus
$\|A\|_F/\|B\|_F\to1$ when $\|B\|_F\to\infty$.
This verifies that residualization is not an asymptotic information
bottleneck in the stated graph sequence.

\subsection*{A rigorous lower-bound subexperiment}
Suppose the allowed compact nuisance set contains $\alpha=\beta=\gamma=0$
and a fixed positive $\sigma$. In that subexperiment the conditional mean
is zero for every $\delta$. In the eigenbasis, the variance ratio relative
to the null is $v_i=(1-\delta b_i)^{-2}$. Gaussian KL divergence is
\[
 \operatorname{KL}(P_\delta,P_0)
 =\tfrac12\sum_i(v_i-1-\log v_i)
 \leq C_\eta\delta^2\|B\|_F^2
\]
for sufficiently small $|\delta|$, by Taylor's theorem uniformly over
$|b_i|\leq1$. Hence if $\delta\|B\|_F\to0$, total variation tends to
zero. A level-$\alpha_0$ test then has power at most $\alpha_0+o(1)$
along those alternatives, by the definition of total variation and
Pinsker's inequality. This establishes a necessary separation order for
any larger model containing the subexperiment. If the fixed compact set
excludes that subexperiment, the lower-bound claim needs adaptation and is
not asserted automatically.

\subsection*{Remaining gap and graph checks}
For disjoint unit-weight edges, eigenvalues are $\pm1$ and
$\|B\|_F^2=N$, giving the parametric candidate order $N^{-1/2}$.
A complete graph normalized by $N-1$ has Frobenius norm squared
$N/(N-1)$, so merely adding vertices does not produce vanishing separation
in the zero-mean subexperiment; it violates the catalogue's divergence
assumption. These checks match the information calculation.

What is not proved is uniform power of the residual test over all separated
alternatives and all mean nuisances. Away from the null,
$(I-\delta B)^{-1}D\theta$ need not lie in the null mean span, and its
quadratic contribution must be bounded jointly with changing covariance.
An efficient score calculation does not justify skipping that step.
Likelihood-ratio profiling with a uniform curvature and stochastic-error
bound is the next concrete route. Heterogeneous and non-Gaussian mechanisms
require additional identifiable innovation restrictions.
\textbf{Final status: UNRESOLVED.}

\section{Completion Estimate}
48\%

This is a post hoc, heuristic estimate for the full catalogue target.
The attempt establishes covariance identification, efficient null information, an exact-size nuisance-free conditional test, and a rigorous separation lower bound in a permitted zero-mean subexperiment. Uniform power over separated alternatives and all mean nuisances, and identifiable heterogeneous or non-Gaussian extensions, are not proved.

\subsection*{Source}
S. Bhadra and M. Schweinberger, \emph{Causal Inference Under Network
Interference}, arXiv:2508.06808v2 (2026), Section 8.1 and the preceding
comparison of contagion with treatment spillover.
\url{https://arxiv.org/html/2508.06808v2#S8.SS1}.

\end{document}
